\documentclass[11pt]{article}
\usepackage{preamble}
\setlength{\parskip}{4pt}

\begin{document}

\daybanner{AP Statistics \textperiodcentered\ Unit 2 \textperiodcentered\ Topic 2.6 \textperiodcentered\ B: 2026-10-26 \textperiodcentered\ E: 2026-11-13 \textperiodcentered\ TEACHER KEY}{Which Way Does the Condition Point?}

\sectionbanner{CED TETHER}

{\small
\begin{itemize}[leftmargin=1.5em,itemsep=2pt,topsep=2pt]
  \item \textbf{Skill 3.A} --- Determine relative frequencies, proportions, or probabilities using simulation or calculations.
  \item \textbf{Enduring Understanding VAR-4} --- The likelihood of a random event can be quantified.
  \item \textbf{LO VAR-4.D} --- Calculate conditional probabilities. [Skill 3.A]
  \item \textbf{EK VAR-4.D.1} --- The probability that event A will occur given that event B has occurred is called a conditional probability and denoted P(A | B) = P(A $\cap$ B) / P(B).
  \item \textbf{EK VAR-4.D.2} --- The multiplication rule states that the probability that events A and B both will occur is equal to the probability that event A will occur multiplied by the probability that event B will occur, given that A has occurred. This is denoted P(A $\cap$ B) = P(A) $\cdot$ P(B | A).
\end{itemize}
}
\textbf{Essential question:} Why can reversing the condition produce a very different probability, even when both probabilities use the same table?\par
\textbf{DOK-3 skill:} 4.B \textperiodcentered\ \textbf{FRQ pattern:} {\small\texttt{conditional-\allowbreak{}direction-\allowbreak{}confusion-\allowbreak{}p-\allowbreak{}a-\allowbreak{}given-\allowbreak{}b-\allowbreak{}vs-\allowbreak{}b-\allowbreak{}given-\allowbreak{}a}}\par
\textbf{DOK rationale:} Part (c) weighs competing claims using the day's numerical or graphical evidence and explains a limitation or consequence; the judgment requires linked reasoning beyond a calculation.\par

\frameworkphaseheader{Do Now (first take) $\rightarrow$ Explore (video + follow-along) $\rightarrow$ Exit (finish + turn in)}{1 $\rightarrow$ 3}{45}{%
\begin{itemize}[leftmargin=1.5em,itemsep=2pt,topsep=2pt]
  \item Show the board at the bell and collect a one-sentence first take without confirming a claim.
  \item Run the named video follow-along, then allow about 10 minutes for parts (a)--(c).
  \item Ask for evidence linked to a claim. Collect the sheet and hand-score part (c) with E/P/I.
\end{itemize}
}{%
\begin{itemize}[leftmargin=1.5em,itemsep=2pt,topsep=2pt]
  \item Choose a claim and cite one cell or margin from the screening table.
  \item Complete the follow-along about conditional probability and the general multiplication rule.
  \item Finish (a)--(c), revise the first take in the exit line, and turn in the sheet.
\end{itemize}
}{%
\begin{itemize}[leftmargin=1.5em,itemsep=2pt,topsep=2pt]
  \item Which number or visible feature supports your claim?
  \item Why does the other claim go beyond that evidence?
  \item What would you tell someone who chose the other claim?
\end{itemize}
}{Circulate and ask students to connect their choice to evidence. Score part (c) using the three required reasoning elements; do not require dropped extension work.}

\begin{callout}{\IconWarn\ WATCH FOR}{calloutred}
\begin{itemize}[leftmargin=1.5em,itemsep=2pt,topsep=2pt]
  \item A choice with copied numbers but no explanation of how they support it.
  \item Treating a model or average as a guaranteed individual outcome.
\end{itemize}
\end{callout}

\sectionbanner{SCORING PART (c) --- E / P / I}

\textbf{Expected elements}
\begin{itemize}[leftmargin=1.5em,itemsep=2pt,topsep=2pt]
  \item \textbf{reasoning-1} (required) --- Rejects Mina using 45/140 or about 32.1 percent
  \item \textbf{reasoning-2} (required) --- Uses 45 with the condition plus 95 without to form the positive group
  \item \textbf{reasoning-3} (required) --- Explains the change in conditioning direction and denominator
\end{itemize}

{\small\renewcommand{\arraystretch}{1.25}
\noindent\begin{tabular}{|p{0.5in}|p{5.9in}|}
\hline
\textbf{E} & Includes all three required reasoning elements above, with correct contextual evidence. \\ \hline
\textbf{P} & Includes at least one substantive correct reasoning link above, but omits or misstates another required element. \\ \hline
\textbf{I} & Includes no substantive correct reasoning link above; an unsupported choice or copied number alone is insufficient. \\ \hline
\end{tabular}}

\textbf{Common mistakes}
\begin{itemize}[leftmargin=1.5em,itemsep=2pt,topsep=2pt]
  \item Using the sensitivity $45/50$ as the chance of condition after a positive result
  \item Dividing 45 by all 1,000 screened students instead of by the 140 positive tests
  \item Calling the 95 positives among students without the condition true positives
  \item Ignoring the 5 percent prevalence when explaining why the two conditional probabilities differ
\end{itemize}

\clearpage
\focusbanner{TODAY'S PROBLEM --- annotated}

Suppose a school screens 1{,}000 students for a hypothetical condition. Exactly $5\%$ have the condition. The test has $90\%$ sensitivity: among students with the condition, $90\%$ test positive. Its false-positive rate is $10\%$: among students without the condition, $10\%$ test positive. The implied counts are shown.\par\smallskip \textbf{Mina:} ``The test catches $90\%$ of students who have the condition, so a positive result means a student has a $90\%$ chance of having it.''\par \textbf{Luis:} ``We need to count everyone who tested positive before finding the chance of having the condition.''\par
\begin{center}
\textbf{Screening results (counts)}\par\smallskip
\begin{tabular}{|l|c|c|c|}
\hline
\textbf{Status} & \textbf{Pos.} & \textbf{Neg.} & \textbf{Total} \\ \hline
\textbf{Condition} & 45 & 5 & 50 \\ \hline
\textbf{No condition} & 95 & 855 & 950 \\ \hline
\textbf{Total} & 140 & 860 & 1{,}000 \\ \hline
\end{tabular}
\end{center}
\begin{firsttakebox}{FIRST TAKE (5 min)}
Does a positive result mean a $90\%$ chance of having the condition? Choose a claim and cite one table count.\par
\answer{Not graded. The familiar 90 percent number is intentionally tempting; the revision should use the full positive-test group.}
\end{firsttakebox}

\textit{Rules callout ``THE GIVEN EVENT SETS THE DENOMINATOR (keep on the page)'' is printed on the student sheet.}\par\medskip

\dokbadge{1}~\textbf{(a)} For $P(\text{condition}\mid\text{positive})$, identify the numerator and denominator counts.\par
\answer{The numerator is 45 students who both have the condition and tested positive. The denominator is all $45+95=140$ students who tested positive.}

\dokbadge{2}~\textbf{(b)} Find $P(\text{positive}\mid\text{condition})$ and $P(\text{condition}\mid\text{positive})$. Show each count fraction and percent.\par
\answer{$P(\text{positive}\mid\text{condition})=45/50=0.90=90\%$. In the other direction, $P(\text{condition}\mid\text{positive})=45/140=9/28\approx0.321=32.1\%$.}

\dokbadge{3}~\textbf{(c)} Is Mina's claim supported? Use the positive-test counts for students with and without the condition. Explain how reversing given changes the group in the denominator. Write 3--4 sentences.\par
\answer{Luis is right: Mina's $90\%$ uses the 50 students with the condition as the denominator. After a positive result, the relevant group is the 140 students who tested positive: 45 have the condition and 95 do not. Thus the chance of the condition given positive is $45/140\approx32.1\%$, not $90\%$. Reversing the given event changes the denominator and answers a different question.}

\begin{summaryexitbox}{EXIT}
One thing the video changed about my first take: \hrulefill
\end{summaryexitbox}

\end{document}
